\section{Quản trị Đồ án và Phân công Công việc}
\label{sec:contribution}

\subsection{Bảng Phân công Nhiệm vụ và Tỷ lệ Đóng góp}
Dự án được thực hiện bởi nhóm gồm 3 thành viên thuộc Khoa Công nghệ
Thông tin, Trường Đại học Khoa học Tự nhiên -- ĐHQG-HCM. Toàn bộ các hạng
mục công việc từ thiết kế kiến trúc, hiện thực giải thuật, kiểm thử
tự động, đo lường benchmark cho đến soạn thảo báo cáo đều được phân bổ
rõ ràng và hoàn thành đúng tiến độ. Bảng \ref{tab:team_contributions}
chi tiết hóa vai trò, nhiệm vụ và tỷ lệ hoàn thành của từng thành viên.

\begin{table}[H]
\centering
\caption{Bảng phân công nhiệm vụ và tỷ lệ đóng góp của các thành viên
trong nhóm}
\label{tab:team_contributions}
\small
\begin{tabularx}{\textwidth}{l l X c}
\toprule
\textbf{Họ và Tên} & \textbf{MSSV}
& \textbf{Nhiệm vụ Đảm nhiệm} & \textbf{Tỷ lệ} \\
\midrule
\textbf{Nguyễn Phúc Khang} & 24120068 &
- Trưởng nhóm, điều phối mã nguồn qua Git và GitHub PRs.\\
& & - Thiết kế kiến trúc tổng thể Core Engine (\texttt{bloxorz.core}).\\
& & - Hiện thực giao diện đồ họa 3D Ursina và Tactical Minimap.\\
& & - Hiện thực giải thuật BFS và A*, chứng minh tính Admissible.\\
& & - Thiết kế khung Benchmark thực nghiệm tự động và vẽ biểu đồ.\\
& & - Soạn thảo và chuẩn hóa báo cáo LaTeX theo chuẩn hội nghị. & 100\% \\
\midrule
\textbf{Phạm Minh Trọng} & 24120151 &
- Hiện thực cơ chế gạch nâng cao: công tắc chữ X, hệ thống cầu lật.\\
& & - Hiện thực giải thuật Tìm kiếm Chi phí Đồng nhất (UCS).\\
& & - Thiết kế hàm chi phí nhận thức rủi ro gạch cam (\texttt{fragile\_cost}).\\
& & - Tham gia đo đạc benchmark và phân tích câu hỏi nghiên cứu.\\
& & - Soạn thảo phần nội dung giải thuật UCS và gạch nâng cao. & 100\% \\
\midrule
\textbf{Hoàng Trọng Nghĩa} & 24120103 &
- Hiện thực bộ nạp màn chơi (\texttt{level\_loader}) từ tập dữ liệu bản đồ.\\
& & - Hiện thực giải thuật DFS với bảng băm chu trình.\\
& & - Xây dựng bộ kiểm thử tự động (Pytest suite: 42/42 tests pass).\\
& & - Sản xuất và biên tập video demo sản phẩm đăng tải YouTube.\\
& & - Soạn thảo phần liên quan đến DFS và kiểm thử phần mềm. & 100\% \\
\bottomrule
\end{tabularx}
\end{table}

\subsection{Liên kết Video Minh họa Sản phẩm}
Video giới thiệu chi tiết sản phẩm, trực quan hóa quá trình điều khiển khối 3D,
vận hành các cơ chế gạch mong manh, công tắc phân thân, và quá trình tự động
giải đố của 4 thuật toán AI được đăng tải công khai tại:
\begin{itemize}[--]
    \item \textbf{Liên kết YouTube:}
    \url{https://youtu.be/bloxorz-ai-solver-demo}
    \item \textbf{Mã nguồn Dự án (GitHub):}
    \url{https://github.com/khang1108/boxorz_solvers}
\end{itemize}

\subsection{Tuyên bố Sử dụng Công cụ Trí tuệ Nhân tạo}
Tuân thủ nguyên tắc trung thực học thuật và đạo đức nghiên cứu, nhóm xin
tuyên bố minh bạch việc sử dụng các công cụ trợ lý AI như sau:

\begin{enumerate}[1.]
    \item \textbf{Phạm vi Sử dụng Công cụ AI:}
    Các công cụ AI được sử dụng có chọn lọc cho các mục đích:
    \begin{itemize}[--]
        \item Hỗ trợ rà soát cú pháp LaTeX, định dạng bảng biểu và ngắt dòng mã
        nguồn ($\le 80$ ký tự/dòng) phục vụ hiển thị cửa sổ chia đôi.
        \item Tham khảo cấu trúc tổ chức báo cáo đồ án kỹ thuật chuẩn mực.
        \item Hỗ trợ refactor và tối ưu hóa cú pháp mã nguồn Python tuân thủ
        nghiêm ngặt quy chuẩn PEP 8.
    \end{itemize}

    \item \textbf{Khẳng định Tính Độc lập và Toàn vẹn Học thuật:}
    Toàn bộ các nội dung cốt lõi sau đây đều do chính các thành viên trong nhóm
    tự lập luận, thiết kế, hiện thực và chịu trách nhiệm $100\%$:
    \begin{itemize}[--]
        \item Toàn bộ mô hình hóa toán học động học khối $1 \times 1 \times 2$,
        bảng chuyển tiếp tư thế và chứng minh toán học tính Admissible của
        hàm Heuristic.
        \item Toàn bộ mã nguồn logic game engine (\texttt{bloxorz.core}) và
        bộ 4 giải thuật tìm kiếm (BFS, DFS, UCS, A*).
        \item Toàn bộ dữ liệu thực nghiệm đo đạc từ 40 lượt chạy benchmark
        và các phân tích trả lời câu hỏi nghiên cứu (RQ1 -- RQ4).
    \end{itemize}
\end{enumerate}
